%% ============================================================================
%%  subdocumento-ejemplo/contenido.tex
%%
%%  Ejemplo minimo de un subdocumento autocontenido. Sirve para verificar que
%%  el mecanismo funciona y como molde para los subdocumentos reales.
%%
%%  Reglas de la carpeta:
%%   - Este archivo SIEMPRE se llama contenido.tex y empieza por \chapter.
%%   - Las figuras van en figuras/ y se llaman por ruta relativa.
%%   - Los fragmentos van en tablas/ o donde quieras, y se traen con \parte{}.
%%   - NADA de aqui conoce la ruta de montaje: no escribas rutas que suban
%%     con ../ ni que empiecen desde la raiz del repositorio. Si lo haces,
%%     mover la carpeta rompe el capitulo.
%% ============================================================================

\chapter{Subdocumento de ejemplo}
\label{cap:ejemplo}

Este capítulo existe para comprobar que el mecanismo de carpetas funciona.
Al montar el informe real, borra esta carpeta y la línea que la llama en
\texttt{main.tex}.

\nota{La carpeta se llama desde \texttt{main.tex} con una sola línea:
\texttt{\textbackslash subdocumento\{subdocumento-ejemplo\}}. Para reubicarla,
mueve la carpeta y cambia esa ruta. No hay nada más que tocar.}

\section{Fragmentos dentro del mismo subdocumento}

Cuando un capítulo crece —un catálogo de requerimientos, una EDT larga— se
parte en archivos dentro de su propia carpeta y se traen con
\texttt{\textbackslash parte\{\}}, que resuelve la ruta \textbf{relativa a esta
carpeta} y no a la de \texttt{main.tex}:

\parte{tablas/ejemplo_tabla}

\section{Figuras}

Las imágenes van en \texttt{figuras/} y se invocan por ruta relativa, porque
\texttt{\textbackslash subdocumento} agrega esa carpeta al
\texttt{\textbackslash graphicspath} mientras dura la importación. Dos
subdocumentos pueden tener cada uno su \texttt{figuras/diagrama.png} sin
pisarse.
